Ethereum's EIP-7702 (activated with the Pectra upgrade on May 7, 2025) fundamentally alters the security model of externally owned accounts (EOAs) by allowing them to delegate code execution to smart contracts via a new transaction type (\texttt{SET\_CODE\_TX\_TYPE}, 0x04). This breaks foundational assumptions that existing Solidity code relies upon: that EOAs cannot execute code, that \texttt{tx.origin == msg.sender} guarantees a non-contract caller, that \texttt{extcodesize} distinguishes EOAs from contracts, and that ETH transfers to EOAs cannot trigger callbacks.

We present ChainEDR, a static analyzer purpose-built for EIP-7702 vulnerability detection. We define a \textit{sandbox boundary model} that captures eight broken assumption classes as systematic boundaries, and define 30 concrete vulnerability checks---21 for EIP-7702 delegation patterns and 9 for ERC-7562 validation scope rules---each mapping a pre-7702 assumption to a post-7702 exploitable pattern with CWE classification and a \textit{PreBehaviorTrace} evidence model.

We construct \textit{EIP7702-Bench}, a labeled benchmark corpus of 47 Solidity contracts (21 vulnerable, 21 fixed, 5 benign). On this corpus, ChainEDR achieves precision\,=\,1.00, recall\,=\,1.00, F1\,=\,1.00 in target-rule evaluation mode over 21 benchmarked checks (Wilson 95\% CI: [0.845,\,1.000]). On production code (499 files, 60,407 LOC across OpenZeppelin Contracts v5.x, MetaMask Delegation Framework, ERC-4337 reference, and Alchemy Modular Account V2), ChainEDR identifies 272 EIP-7702-specific findings requiring manual triage, including 4 CRITICAL-severity patterns verified through manual triage. The compared tools---Slither, Aderyn, Mythril, and Semgrep---do not ship dedicated EIP-7702 detectors and report no EIP-7702-specific findings in our configuration. ChainEDR integrates with CI/CD pipelines via SARIF v2.1.0 output.
